\section{Related Work}

Our study connects three research streams: qualitative analyses of benchmark construction, quantitative reproducibility investigations, and concentration measurement in science. Each provides essential background but none addresses whether benchmark concentration creates temporal feedback loops.

\subsection{Benchmark Analysis and Construction}

Sociological and ethnographic studies have examined how benchmarks become established standards. \citet{engdahl2024benchmark} provides a detailed qualitative account of benchmark construction in machine learning, documenting the community processes through which particular datasets become accepted evaluation standards. This work illuminates the social dynamics behind standardization but does not quantify whether these dynamics create lock-in over time.

Technical work on benchmark design addresses quality and coverage. Holistic evaluation platforms such as Dynaboard \citep{ma2021dynaboard} aim to provide richer assessment than single-metric leaderboards, while large-scale benchmark collections like HPO-B \citep{pineda2021hpob} aggregate datasets from sources like OpenML to enable broader hyperparameter optimization research.

\subsection{Reproducibility and Replication Studies}

A substantial literature examines whether machine learning results replicate. \citet{olszewski2023reproducibility} investigate the relationship between reproducibility and artifact evaluation committees, finding limited effects of formal review processes on subsequent replication success. This work tracks methodological practices but focuses on code and experiment availability rather than dataset selection dynamics.

\subsection{Concentration Metrics in Science Studies}

Science of science research has developed measures of concentration in research portfolios, citation patterns, and topic distributions. The Herfindahl-Hirschman Index (HHI), originally from industrial organization economics, has been applied to publication concentration across topics. We adapt HHI to benchmark selection, computing concentration across datasets rather than across topics or authors.

\subsection{The Gap: No Temporal Causality Test}

Existing work establishes that benchmark concentration exists, that benchmarks spread through community processes, and that citation patterns reveal intellectual influence. What no prior study provides is a quantitative test of whether concentration in one time period causally affects diversity in subsequent periods---the defining characteristic of lock-in.
